\documentclass[10pt,a4paper]{amsart}
\usepackage[margin=19mm]{geometry}
\usepackage{amsmath,amssymb,mathtools}
\usepackage[scr=rsfs,cal=euler]{mathalfa}
\usepackage{xfrac}
\usepackage[unicode,hidelinks,bookmarks=false]{hyperref}
\newcommand{\ie}{i.e.\ }
\newcommand{\cf}{cf.\ }
\newcommand{\resp}{respectively\ }
\newcommand{\Subsection}{\subsection}
\providecommand{\nobreakdash}{\nobreak\hbox{-}}
\setlength{\emergencystretch}{2em}
\multlinegap0pt
\usepackage{bm}
\usepackage{algorithm}\usepackage{algpseudocode}
\usepackage{mathtools}
\usepackage[shortlabels]{enumitem}
\usepackage{amssymb}
\newtheorem{proposition}{Proposition}
\DeclareMathOperator{\id}{Id}
\newcommand{\norm}[1]{\|#1\|}
\title{WinQ: Accelerating Quantization-Aware Training of Language Models Around Saddle Points}
\author{Dongyue Li, Zechun Liu, Kai Yi, Zhenshuo Zhang, Changsheng Zhao, Raghuraman Krishnamoorthi, Harshit Khaitan, Hongyang R. Zhang, Steven Li}
\date{Source: arXiv:2605.17471v1 (CC BY 4.0)}
\begin{document}
\maketitle

\noindent\emph{Source and licence.} WinQ: Accelerating Quantization-Aware Training of Language Models Around Saddle Points. Version arXiv:2605.17471v1 (CC BY 4.0), distributed by arXiv under the Creative Commons Attribution 4.0 International licence, http://creativecommons.org/licenses/by/4.0/, as recorded in the arXiv licence metadata for this exact version and shown on the abstract page https://arxiv.org/abs/2605.17471v1. Full licence notice: \texttt{licenses/arxiv-2605.17471-CC-BY-4.0.txt}.\par\smallskip

\noindent\textit{Editorial scope.} Counted unit: the article's proposition prop\_connect (source0.tex lines 878--886). The article's complete original proof of it is delivered with it (source0.tex lines 1220--1254, one \texttt{proof} environment of 35 lines, which the article places away from the statement). No mathematics is written, repaired or re-derived: the statement and the proof are the article's own lines, in the article's own notation, and no retained line is substituted or re-flowed.  No further source-local context is needed: every cross-reference of every retained line resolves inside the two delivered spans.  Nothing was omitted from any retained span, so the package carries no omission record.  No retained line contains a citation, so the wrapper carries no bibliography and no bibliography entry.  No retained line contains a figure environment, an image-inclusion command or drawing code, so no image is delivered and the package declares no asset dependency.  Editorial formatting only, in addition: an AMS \texttt{amsart} preamble that restates the article's own declarations and macros used by the retained lines, namely the environments \texttt{proposition} on separate counters, together with the standard packages those lines need.  The retained environments print in this document as Proposition 1 in order of appearance, whereas the article prints the same items with the numbering of its own full document.  The complete native source of the article as distributed in the arXiv e-print for this exact version is bundled unchanged as \texttt{sources/200797/source0.tex}.
\begin{proposition}\label{prop_connect}
Let $W$ denote the continuous latent weights, $Q(W)$ the corresponding quantized weights, and $\eta > 0$ the learning rate of the optimizer.
Let $q = Q(W)$ denote the quantized weights of weight $W$. Assume that the quantized weights are locally stable, meaning that there exists a neighborhood $\mathcal{N}_q \subseteq \mathbb{R}^d$ of $W$, such that $Q(Z)=q$ for all $Z \in \mathcal{N}_q$.
The weight interpolation step $W \leftarrow (1-\alpha)W + \alpha Q(W)$ with $\alpha \in [0, 1)$ is equivalent to gradient update on the $\ell_2$-regularized objective:
\begin{align}
    \Phi(W) = L_Q(W) + \frac{\gamma}{2} \|W - q\|^2,
\end{align}
where $\alpha$ connects to the regularization strength $\gamma$ as $\alpha = \frac{\eta \gamma}{1 + \eta \gamma}$. Correspondingly, the Hessian of $\Phi(W)$ becomes \[ \nabla^2 \Phi(W) = \nabla^2 L_Q(W) + \gamma \id. \]
\end{proposition}

\begin{proof}[Proof of Proposition \ref{prop_connect}]
We consider optimizing the objective $\Phi(W)$ using proximal gradient descent. We decompose the objective into the loss component $L_Q(W)$, which is optimized via the Straight-Through Estimator, and the regularization component $R(W) = \frac{\gamma}{2} \norm{W - q}_2^2$, where $q = Q(W)$ is the target quantization grid point. We assume $q$ is locally constant within the current quantization cell.

A proximal gradient descent update at iteration $t$ consists of two components: 
\begin{itemize}
    \item Gradient update step: we conduct a gradient step on the loss $L_Q(W)$ using the learning rate $\eta$: 
    \[V = W_t - \eta \nabla L_Q(W_t).\]
    \item Proximal update step: we apply the proximal operator of the scaled regularizer $\eta R(W)$ to the intermediate weight $V$:
    \[W_{t+1} = \text{prox}_{\eta R}(V) = \arg\min_W \left( R(W) + \frac{1}{2\eta} \|W - V\|_2^2 \right).\]
\end{itemize} 

Substituting the $\ell_2$ penalty $R(W)$ into the proximal operator yields:
\begin{align*}
W_{t+1} = \arg\min_W \left( \frac{\gamma}{2} \norm{W - q}_2^2 + \frac{1}{2\eta} \norm{W - V}_2^2 \right).    
\end{align*}
As the objective is a convex quadratic function, this has a unique global minimum where 
\begin{align*}
\nabla_W \left( \frac{\gamma}{2} \|W - q\|_2^2 + \frac{1}{2\eta} \|W - V\|_2^2 \right) = \gamma(W - q) + \frac{1}{\eta}(W - V) = 0    
\end{align*}
Therefore, we have
\begin{align*}
W = \left( \frac{1}{1 + \eta \gamma} \right) V + \left( \frac{\eta \gamma}{1 + \eta \gamma} \right) q. 
\end{align*}

When setting the interpolation scalar $\alpha$ as $\frac{\eta \gamma}{1 + \eta \gamma}$. Substituting $\alpha$ and $1 - \alpha$ back into our equation yields the exact weight interpolation step:
\begin{align*}
W_{t+1} = (1 - \alpha) V + \alpha q.     
\end{align*}
This shows that the weight interpolation is mathematically equivalent to a proximal update applied to the $\ell_2$ penalty.
Additionally, the second derivative of the quadratic penalty term is $\gamma I$, which yields
\begin{align*}
\nabla^2 \Phi(W) = \nabla^2 L_Q(W) + \gamma \id.     
\end{align*}
This, in turn, increases the magnitude of the Hessian eigenvalues.
\end{proof}

\end{document}
